\begin{abstract}

% Smart home platforms rely on automation rules to coordinate devices and sensors, but independently configured automation rules can interact through shared entities and zone-dependent physical-channel propagation, leading to unintended and unsafe behaviors. Existing rule conflict detectors often over-approximate potential interactions without considering zone-specific propagation and user-specific safety preferences. Moreover, they rarely confirm whether a suspected conflict is actually activated at runtime, leading to false positives. We present a hybrid framework that combines offline static analysis with runtime dynamic monitoring for rule conflict detection. The framework incorporates zone-dependent physical-channel propagation and a user-defined Entity Safety Configuration to focus on safety-relevant candidate rule conflicts, and uses runtime assertion verification to confirm activated conflicts under the current system state. Implemented on Home Assistant and evaluated in a multi-zone testbed with 17 automation rules, our framework detects all manually confirmed conflict cases and eliminates nine false positives reported by IoTMediator, while keeping runtime verification overhead at the millisecond level.

Smart home platforms utilize automation rules to coordinate devices and sensors. These rules often interact through shared devices and physical channels, such as temperature or illumination, leading to unintended behaviors. Existing rule conflict detection approaches often identify logical overlaps between automation rules without determining whether the underlying interaction is physically relevant in the current deployment or harmful under the target application scenario. As a result, they may over-report benign interactions as rule conflicts, especially when physical propagation is constrained by zone boundaries or when user-specific safety requirements differ across scenarios. In this paper, we propose a hybrid framework that combines offline static analysis with runtime dynamic monitoring. The framework utilizes a Trigger-Condition-Action-Environment (TCAE) model to account for zone-dependent physical interactions. We introduce an entity safety configuration that allows users to define safety boundaries, transforming potential interactions into a Rule Conflict Dependency Graph (RCDG). A runtime monitor then performs assertion verification to confirm whether suspected conflicts are actually triggered under the current system state. We implemented the framework on Home Assistant and evaluated it in a multi-zone environment. The results show that our approach identifies all confirmed conflicts and eliminates nine false positives reported by the state-of-the-art tool IoTMediator. The runtime verification maintains a millisecond-level overhead, ensuring system efficiency.

\keywords{Smart home · rule conflict · hybrid detection · runtime verification}

\end{abstract}